\subsection{What varying the constant twist does and does not prove}

\begin{claim}{B.TW1}{Proved}{A uniform count of bad constant twists}
Let $C/K$ be smooth proper of genus $g\ge2$ over a finite field
of characteristic $p$. Let $D\to C$ be a geometrically connected
\emph{\'etale} cover with constant cyclic deck group of odd prime
order $\ell\ne p$. Put $f=\operatorname{ord}_\ell(p)$.
Among the $\ell$ constant twists of this cover, at most
\[
                \frac{(g-1)(\ell-1)}{f}
\]
have relative class number divisible by $p$. In particular, if
$f\ge g-1$, at least one twist has $p$-unit relative class number.
For genus two this last condition is automatic.
\end{claim}
\begin{proof}
Fix a generator $\delta$ of the deck group. The nontrivial factor
of its group algebra is
\[
 R=\mathbf F_p[X]/\Phi_\ell(X)
       \simeq\prod_{j=1}^{(\ell-1)/f}\mathbf F_{p^f}.
\]
Let $W$ be the new part of
$H^1_{\et}(D_{\bar K},\mathbf F_p)$. Over an algebraic closure,
the nontrivial coherent character spaces are $H^1(C,L)$ for
nontrivial degree-zero line bundles $L$ of order $\ell$; each
has dimension $g-1$. The Artin--Schreier comparison embeds
\'etale cohomology tensored with $\bar K$ into the Frobenius-stable
coherent part, equivariantly for the deck action. Consequently the rank $d_j$ of $W$ over each field
factor of $R$ satisfies $0\le d_j\le g-1$.

The deck group is constant, so arithmetic Frobenius $F$ commutes
with its action and is linear over every field factor. Twisting
by the constant cocycle with Frobenius value $\delta^t$, for
$t\in\mathbf Z/\ell$, replaces its action on $W$ by
$\delta^tF$, up to using $-t$ consistently in the contravariant
cohomology convention. On the $j$th field factor, write the image
of $\delta$ as the primitive root $\zeta_j$. A new fixed vector
exists on this factor exactly when
\[
               \det(\zeta_j^{-t}-F\mid W_j)=0.
\]
This polynomial has degree $d_j$, and its $\ell$ arguments are
distinct. At most $d_j$ twists are therefore bad on this factor.
Taking the union over the $(\ell-1)/f$ factors gives the claimed
bound. Nonsemisimplicity of $F$ causes no problem: a fixed vector
is detected exactly by this determinant. The rational two-primary
argument in \statementref{B.W1} applies with $p$ in place of two:
norm and pullback split the $p$-primary groups because
$p\nmid\ell$, and rational new $p$-torsion is equivalent to
divisibility of the relative class number. Finally,
$f\ge g-1$ makes the bound at most $\ell-1$, strictly below the
number of twists.
\end{proof}

\begin{claim}{B.TW2}{Proved}{Existence with an unspecified base point}
For every smooth proper curve $C/k$ of genus $g\ge2$ over a
finite field, there are a finite extension $K/k$, a point
$P_K\in C(K)$, and a constant cyclic $P_K$-split cover of odd
prime degree, different from the characteristic, with $p$-unit
relative class number. The point $P_K$ is not prescribed in
advance. More quantitatively, for a constant degree-$\ell$ cover
as in \statementref{B.TW1}, put
$B=(g-1)(\ell-1)/f$ and $g_D=1+\ell(g-1)$. If $f\ge g-1$, the
number of $P\in C(K)$ for which its unique split twist is a
class-number witness is at least
\[
       \frac{\ell-B}{\ell}
             \bigl(\#K+1-2g_D\sqrt{\#K}\bigr).
\]
\end{claim}
\begin{proof}
For each twist $D_t$, let $N_t$ be the number of rational points
of $C$ that split in it. Every geometric fibre has $\ell$ points
and Frobenius is a deck translation. Thus each $P\in C(K)$
splits in precisely one twist, and
$\#D_t(K)=\ell N_t$. At least $\ell-B$ twists have $p$-unit
relative class number. Applying the Weil lower bound to each
of these geometrically connected curves gives the displayed
inequality \cite[Theorem 1.6 and formula 1.5.1]{WeilCurves}.
If its right-hand side is negative, the inequality is automatic;
otherwise summing over at least $\ell-B$ good twists proves it.

There exist odd primes $\ell\ne p$ with
$\operatorname{ord}_\ell(p)\ge g-1$: the primes with order less
than $g-1$ divide the fixed nonzero integer
$\prod_{j=1}^{g-2}(p^j-1)$ and are finite in number.
Choose such an $\ell$ and a nonzero geometric cyclic
$\ell$-cover, which exists because the prime-to-$p$ abelian
fundamental group has rank $2g$ at $\ell$.
It, its deck group, and a model of its base curve descend to a
finite extension. Enlarge that field until
$\#K+1>2g_D\sqrt{\#K}$. The quantitative assertion then produces
a good twist with a rational point above some $P_K$.
All descent and field enlargements precede the choice of a good
twist; no assertion that a previously good twist remains good
under extension is needed.
\end{proof}

\paragraph{The remaining pointing obstruction in B1.}
For a fixed prescribed $P$, its split twist is one specified
element of the $\ell$ twists. The counting theorem can leave
that element among the bad twists. There is no way to repair
this by simply extending a good unpointed twist until $P$
splits: this extension may make its nontrivial $p$-torsion
eigenvalues equal to one.

This failure occurs on the explicit R8 curve, not just on an
abstract representation. Consider the cyclic degree-three
character obtained by pulling back the point $(0,0)$ of exact
order three on $E:y^2+y=x^3$ along its first degree-two map
$C\to E$. Norm followed by pullback is multiplication by two,
so this character is nonzero. Its pointed cover is constant over
$k=\mathbf F_{2^{38}}$, since $s=-2^{19}\equiv1\pmod3$.
Its new geometric two-rank is positive by \statementref{B.R5};
the nontrivial group algebra is the single field $\mathbf F_4$,
and the coherent character dimension is one, so its new
\'etale cohomology has rank one over that field. The commuting
pointed hyperelliptic involution used in \statementref{B.W3}
forces Frobenius to be the identity. Both nontrivial constant
twists instead have primitive third-root Frobenius multiplier,
and hence odd relative class number. They do not split over $P$.
The cubic constant extension first makes that fibre split and
simultaneously cubes the multiplier to one, making the relative
class number even again. Thus the unpointed existence theorem
does not settle B1, including the range $p>(2g-2)^2$.

These three covers can be given explicitly in the original R8
coordinates by
\[
       w^3=y,\qquad w^3=ty,\qquad w^3=t^2y.
\]
On the elliptic curve, $\operatorname{div}(y)=3(0,0)-3O$;
its pullback gives the exact degree-three character above.
If $\tau=x/y$ is the elliptic local parameter at $O$, then
$y\tau^3$ has residue one. The degree-two map has
$\tau=b\eta^2+\cdots$ at $P$ for a rational local parameter
$\eta$ and $b\in k^\times$. The leading unit of $y$ at $P$
is therefore $b^{-3}$, a cube, proving that the first cover
is precisely the split twist.
The pure binary-polynomial program
\path{code/criterion_twist_obstruction.py} verifies irreducibility
of $M$ and computes
\[
        t^{(2^{38}-1)/3}=588964846
\]
in the polynomial-bit encoding of $k$. This value is a nontrivial
cube root of unity, so $t$ is a noncube and the other two covers
first split above $P$ over the cubic extension. The same program
checks the two-dimensional new cohomology matrices
\[
 I,\qquad
 \begin{pmatrix}0&1\\1&1\end{pmatrix},\qquad
 \begin{pmatrix}1&1\\1&0\end{pmatrix},
\]
up to interchanging the last two by the chosen cohomology
convention. Their determinants $\det(1-F)$ are $0,1,1$ in
$\mathbf F_2$, and all three matrices cube to $I$.
The full finite data and the geometric scope of this check are
recorded in its JSON and run transcript.
